% Unit 046 terjemahan Bahasa Indonesia dari chapter3.tex baris 2936--3200.
% Sumber: Methods of Algebra, Volume 2, commit
% 9a5803ff77dd3257484cb177f851a73770a59dd3 (CC BY 4.0).
% stable-unit-id: o014.aljabr2.chapter3.k-injectives
% Cakupan: Bagian 3.15 lengkap.

% segment-id: o014.li2.chapter3.k-injectives.h001
\section{Kompleks K-Injektif dan K-Projektif}\label{sec:K-injectives}
\hypertarget{o014.aljabr2.chapter3.k-injectives}{ }

% segment-id: o014.li2.chapter3.k-injectives.p001
Resolusi injektif dan projektif bagi kompleks memberikan definisi awal
funktor turunan dalam \S\ref{sec:derived-primer}. Konstruksi tersebut
mensyaratkan kompleks terbatas bawah atau terbatas atas, sedangkan penerapan
yang melibatkan kompleks tak terbatas beserta kategori turunannya membutuhkan
kelas resolusi yang lebih luas. Teori resolusi K-injektif dan K-projektif
bermula dari \cite{Spa88}. Bagian ini bertujuan memberikan sejumlah
persiapan terkait; argumennya mengikuti \cite{stacks}, dan pembaca juga dapat
merujuk \cite[Chapter 14]{KS06} atau \cite{BN93}.

% segment-id: o014.li2.chapter3.k-injectives.de001
\begin{definisi}[N.\ Spaltenstein]\label{def:K-resolutions}
	\index{kompleks!K-injektif, K-projektif (K-injective, K-projective)}
	\index{resolusi!K-injektif, K-projektif (K-injective, K-projective)}
	\index{cukup banyak kompleks K-injektif (enough K-injectives)}
	\index{cukup banyak kompleks K-projektif (enough K-projectives)}
	Misalkan $\mathcal{A}$ kategori abelian.
% segment-id: o014.li2.chapter3.k-injectives.l001
	\begin{itemize}
% segment-id: o014.li2.chapter3.k-injectives.i001
		\item Kompleks $I \in \Obj(\cate{C}(\mathcal{A}))$ disebut
		\emph{kompleks K-injektif}\footnote{Istilah yang mungkin lebih wajar
		adalah kompleks injektif homotopik, dan secara serupa kompleks
		projektif homotopik.} jika, untuk setiap kompleks asiklik
		$X \in \Obj(\cate{C}(\mathcal{A}))$, kompleks
		$\Hom^\bullet(X,I)$ juga asiklik.

		Jika $f:A\to I$ merupakan kuasi-isomorfisme di
		$\cate{C}(\mathcal{A})$ dan $I$ K-injektif, maka $f$ disebut
		\emph{resolusi K-injektif} dari $A$.
% segment-id: o014.li2.chapter3.k-injectives.i002
		\item Kompleks $P \in \Obj(\cate{C}(\mathcal{A}))$ disebut
		\emph{kompleks K-projektif} jika, untuk setiap kompleks asiklik
		$X \in \Obj(\cate{C}(\mathcal{A}))$, kompleks
		$\Hom^\bullet(P,X)$ juga asiklik.

		Jika $f:P\to A$ merupakan kuasi-isomorfisme di
		$\cate{C}(\mathcal{A})$ dan $P$ K-projektif, maka $f$ disebut
		\emph{resolusi K-projektif} dari $A$.
	\end{itemize}

	Jika setiap objek $\cate{C}(\mathcal{A})$ mempunyai resolusi K-injektif
	(atau resolusi K-projektif), maka $\mathcal{A}$ dikatakan mempunyai
	\emph{cukup banyak kompleks K-injektif} (atau
	\emph{cukup banyak kompleks K-projektif}).
\end{definisi}

% segment-id: o014.li2.chapter3.k-injectives.p002
Kedua rangkaian definisi itu tentu saling dual. Dalam praktik, kriteria
berikut sering kali lebih mudah digunakan.

% segment-id: o014.li2.chapter3.k-injectives.le001
\begin{lema}\label{prop:K-injective-criterion}
	Kompleks $I$ (atau $P$) bersifat K-injektif (atau K-projektif) jika dan
	hanya jika, untuk setiap kompleks asiklik $X$,
	$\Hom_{\cate{K}(\mathcal{A})}(X,I)=0$ (atau
	$\Hom_{\cate{K}(\mathcal{A})}(P,X)=0$).
\end{lema}

% segment-id: o014.li2.chapter3.k-injectives.pf001
\begin{bukti}
	Misalkan $n\in\Z$. Lema \ref{prop:Hom-cplx-d-translate} memberikan
	\[
	\Hm^n \Hom^\bullet(X, I) \simeq
	\Hom_{\cate{K}(\mathcal{A})}(X[-n], I), \quad
	\Hm^n \Hom^\bullet(P, X) \simeq
	\Hom_{\cate{K}(\mathcal{A})}(P, X[n]);
	\]
	akan tetapi, $X$ asiklik jika dan hanya jika $X[-n]$ asiklik, dan ini
	juga setara dengan $X[n]$ asiklik.
\end{bukti}

% segment-id: o014.li2.chapter3.k-injectives.ex001
\begin{contoh}\label{eg:bdd-K-resolution}
	Lema \ref{prop:resolution-homotopy-aux} langsung dapat dirumuskan ulang
	sebagai berikut.
% segment-id: o014.li2.chapter3.k-injectives.l002
	\begin{compactitem}
% segment-id: o014.li2.chapter3.k-injectives.i003
		\item Jika $I\in\Obj(\cate{C}^+(\mathcal{A}))$ tersusun atas
		objek-objek injektif, maka $I$ K-injektif.
% segment-id: o014.li2.chapter3.k-injectives.i004
		\item Jika $P\in\Obj(\cate{C}^-(\mathcal{A}))$ tersusun atas
		objek-objek projektif, maka $P$ K-projektif.
	\end{compactitem}
	Dengan demikian, resolusi injektif (atau projektif) merupakan kasus khusus
	resolusi K-injektif (atau K-projektif).
\end{contoh}

% segment-id: o014.li2.chapter3.k-injectives.ex002
\begin{contoh}[A.\ Dold]
	Syarat terbatas pada satu sisi dalam Contoh
	\ref{eg:bdd-K-resolution} sangat diperlukan. Sebagai contoh, ambil
	$\mathcal{A}=\Z/4\Z\dcate{Mod}$ dan tinjau kompleks asiklik $Q$
	di dalamnya:
	\[
	\cdots \xrightarrow{2} \Z/4\Z \xrightarrow{2} \Z/4\Z
	\xrightarrow{2} \Z/4\Z \xrightarrow{2} \cdots.
	\]
	Setiap sukunya merupakan modul bebas. Di sisi lain, hasil kali tensor
	suku demi suku $Q \dotimes{\Z/4\Z} \Z/2\Z$ memberikan kompleks
	\[
	\cdots \xrightarrow{0} \Z/2\Z \xrightarrow{0} \Z/2\Z
	\xrightarrow{0} \Z/2\Z \xrightarrow{0} \cdots.
	\]
	Oleh sebab itu, $Q$ tidak K-projektif. Jika $Q$ K-projektif, maka
	$\identity_Q$ homotopik dengan $0$, sehingga
	$\identity_{Q \otimes \Z/2\Z}$ juga demikian, bertentangan dengan
	$\Hm^n(Q \dotimes{\Z/4\Z} \Z/2\Z)=\Z/2\Z$.

	Dari sudut pandang lain, contoh ini juga memperlihatkan bahwa kompleks tak
	terbatas yang tersusun atas objek-objek projektif bukanlah resolusi yang
	layak: $0\to Q$ dan $0\to0$ keduanya merupakan kuasi-isomorfisme, tetapi
	paragraf sebelumnya menunjukkan bahwa $\identity_Q$ tidak homotopik dengan
	$0$. Jadi, $Q$ dan $0$ tidak isomorfik dalam
	$\cate{K}(\mathcal{A})$, dan Teorema \ref{prop:resolution-homotopy}
	tidak dapat begitu saja diperluas ke kasus tak terbatas.
\end{contoh}

% segment-id: o014.li2.chapter3.k-injectives.p003
Hasil berikut merupakan perumuman alami Teorema
\ref{prop:resolution-homotopy}; pembuktiannya juga ditangguhkan hingga
pembahasan setelah Teorema
\sourcecrossref{prop:cplx-triangulated}{pada \S~4.4}.

% segment-id: o014.li2.chapter3.k-injectives.th001
\begin{teorema}\label{prop:K-resolution-homotopy}
	Tetapkan kuasi-isomorfisme $\alpha:X\to Y$ di
	$\cate{C}(\mathcal{A})$.
% segment-id: o014.li2.chapter3.k-injectives.l003
	\begin{itemize}
% segment-id: o014.li2.chapter3.k-injectives.i005
		\item Diberikan morfisme $\gamma:X\to I$, dengan $I$ kompleks
		K-injektif, terdapat diagram komutatif dalam
		$\cate{K}(\mathcal{A})$
% segment-id: o014.li2.chapter3.k-injectives.g001
		\[\begin{tikzcd}[row sep=tiny]
			& Y \arrow[dd, "\beta"] \\
			X \arrow[ru, "\alpha"] \arrow[rd, "\gamma"'] & \\
			& I
		\end{tikzcd}\]
% segment-id: o014.li2.chapter3.k-injectives.i006
		\item Diberikan morfisme $\gamma:P\to Y$, dengan $P$ kompleks
		K-projektif, terdapat diagram komutatif dalam
		$\cate{K}(\mathcal{A})$
% segment-id: o014.li2.chapter3.k-injectives.g002
		\[\begin{tikzcd}[row sep=tiny]
			& X \arrow[ld, "\alpha"'] \\
			Y & \\
			& P \arrow[lu, "\gamma"] \arrow[uu, "\beta"']
		\end{tikzcd}\]
	\end{itemize}
	Dalam kedua kasus, $\beta$ unik sebagai morfisme di
	$\cate{K}(\mathcal{A})$.
\end{teorema}

% segment-id: o014.li2.chapter3.k-injectives.pr001
\begin{proposisi}\label{prop:K-sigma-duality}
	Ekuivalensi $\sigma$ dalam Definisi--Proposisi \ref{def:sigma} memenuhi
	sifat berikut: $X\in\Obj(\cate{C}(\mathcal{A}^{\opp}))$ bersifat
	K-injektif (atau K-projektif) jika dan hanya jika
	$\sigma X\in\Obj(\cate{C}(\mathcal{A})^{\opp})
	=\Obj(\cate{C}(\mathcal{A}))$ bersifat K-projektif (atau K-injektif).
\end{proposisi}

% segment-id: o014.li2.chapter3.k-injectives.pf002
\begin{bukti}
	Funktor $\sigma$ mempertahankan kompleks asiklik dan menginduksi
	$\cate{K}(\mathcal{A}^{\opp})\rightiso
	\cate{K}(\mathcal{A})^{\opp}$ (Proposisi
	\ref{prop:sigma-homotopy}); karena itu, $\sigma$ mempertahankan kriteria
	dalam Lema \ref{prop:K-injective-criterion}.
\end{bukti}

% segment-id: o014.li2.chapter3.k-injectives.p004
Selanjutnya, tinjau funktor-funktor adjoin. Proposisi
\ref{prop:adjoint-injective-projective} mempunyai versi yang bersesuaian
dalam konteks ini.

% segment-id: o014.li2.chapter3.k-injectives.pr002
\begin{proposisi}\label{prop:adjoint-injective-projectives-K}
	Tinjau sepasang funktor di antara kategori-kategori abelian
% segment-id: o014.li2.chapter3.k-injectives.g003
	$\begin{tikzcd}
		\mathcal{A} \arrow[r, shift left, "F"] &
		\mathcal{B} \arrow[l, shift left, "G"]
	\end{tikzcd}$,
	dan andaikan $F$ funktor eksak.
% segment-id: o014.li2.chapter3.k-injectives.l004
	\begin{compactitem}
% segment-id: o014.li2.chapter3.k-injectives.i007
		\item Jika $G$ merupakan adjoin kiri dari $F$, maka
		$\cate{K}G:\cate{K}(\mathcal{B})\to\cate{K}(\mathcal{A})$
		membawa kompleks K-projektif ke kompleks K-projektif;
% segment-id: o014.li2.chapter3.k-injectives.i008
		\item jika $G$ merupakan adjoin kanan dari $F$, maka
		$\cate{K}G:\cate{K}(\mathcal{B})\to\cate{K}(\mathcal{A})$
		membawa kompleks K-injektif ke kompleks K-injektif.
	\end{compactitem}
\end{proposisi}

% segment-id: o014.li2.chapter3.k-injectives.pf003
\begin{bukti}
	Berdasarkan dualitas, cukup tinjau kasus ketika $G$ merupakan adjoin kiri.
	Dalam hal ini, $G$ otomatis merupakan funktor aditif (Korolari
	\ref{prop:automatic-additivity}). Misalkan
	$P\in\Obj(\cate{C}(\mathcal{B}))$ kompleks K-projektif. Untuk setiap
	kompleks asiklik $X\in\Obj(\cate{C}(\mathcal{A}))$, Proposisi
	\ref{prop:KF-adjoint} memberikan isomorfisme
	\[
	\Hom_{\cate{K}(\mathcal{A})}(\cate{K}G(P), X) \simeq
	\Hom_{\cate{K}(\mathcal{B})}(P, \cate{K}F(X)).
	\]
	Karena $F$ eksak, $\cate{K}F(X)$ asiklik, sehingga ruas kanan bernilai
	$0$. Terapkan Lema \ref{prop:K-injective-criterion}.
\end{bukti}

% segment-id: o014.li2.chapter3.k-injectives.p005
Tugas selanjutnya dalam bagian ini adalah menyelidiki keberadaan resolusi
K-injektif atau K-projektif. Argumennya agak berliku; kita mulai dengan kasus
K-injektif.

% segment-id: o014.li2.chapter3.k-injectives.le002
\begin{lema}\label{prop:K-injective-projlimit-aux}
	Tinjau suatu barisan kompleks di atas $\mathcal{A}$,
	$\cdots\to I_2\to I_1\to I_0$. Andaikan:
% segment-id: o014.li2.chapter3.k-injectives.l005
	\begin{itemize}
% segment-id: o014.li2.chapter3.k-injectives.i009
		\item setiap $I_k$ K-injektif;
% segment-id: o014.li2.chapter3.k-injectives.i010
		\item untuk setiap $k$ dan $n$, morfisme
		$I_{k+1}^n\to I_k^n$ merupakan epimorfisme yang mempunyai
		penampang (dengan kata lain, mempunyai invers kanan; lihat pembahasan
		setelah Proposisi \ref{prop:split-ses});
% segment-id: o014.li2.chapter3.k-injectives.i011
		\item untuk setiap $n$, objek
		$I^n:=\varprojlim_k I_k^n$ ada di $\mathcal{A}$.
	\end{itemize}
	Maka diferensial terinduksi
	$d_I^n:=\varprojlim_k d_{I_k}^n:I^n\to I^{n+1}$ menjadikan
	$I:=(I^n,d_I^n)_n$ kompleks K-injektif.
\end{lema}

% segment-id: o014.li2.chapter3.k-injectives.pf004
\begin{bukti}
	Misalkan $X\in\Obj(\cate{C}(\mathcal{A}))$ asiklik. Untuk setiap
	$k\in\Z_{\geq0}$, ambil penggalan berikut dari kompleks
	$\Hom^\bullet\left(X,I_k\right)$:
% segment-id: o014.li2.chapter3.k-injectives.e001
	\begin{equation}\label{eqn:K-injective-projlimit-aux}\begin{tikzcd}
		[row sep=small, column sep=tiny,
		every cell/.append style = {font = \small}]
		\displaystyle\prod_n \Hom\left( X^n, I_k^{n-2} \right) \arrow[r] &
		\displaystyle\prod_n \Hom\left( X^n, I_k^{n-1} \right) \arrow[r] &
		\displaystyle\prod_n \Hom\left( X^n, I_k^n \right) \arrow[r] &
		\displaystyle\prod_n \Hom\left( X^n, I_k^{n+1} \right) \\
		\Hom^{-2}\left( X, I_k \right) \arrow[equal, u] &
		\Hom^{-1}\left( X, I_k \right) \arrow[equal, u] &
		\Hom^0\left( X, I_k \right) \arrow[equal, u] &
		\Hom^1\left( X, I_k \right) \arrow[equal, u]
	\end{tikzcd}\end{equation}
	Barisan ini eksak karena $I_k$ merupakan kompleks K-injektif. Sekarang
	tinjau homomorfisme yang diinduksi oleh $I_{k+1}\to I_k$,
% segment-id: o014.li2.chapter3.k-injectives.q001
	\[
	\Hom\left(X^n,I_{k+1}^{n-2}\right)\to
	\Hom\left(X^n,I_k^{n-2}\right);
	\]
	karena terdapat penampang $I_k^{n-2}\to I_{k+1}^{n-2}$, homomorfisme
	tersebut merupakan epimorfisme. Sifat ini tetap berlaku setelah mengambil
	$\prod_n$. Jadi, ketika $k$ berubah, suku paling kiri dalam
	\eqref{eqn:K-injective-projlimit-aux} memenuhi syarat Mittag--Leffler
	(Definisi \ref{def:ML}).

	Sekarang biarkan $k\in\Z_{\geq0}$ berubah dalam
	\eqref{eqn:K-injective-projlimit-aux}. Korolari
	\ref{prop:ML-4-terms} memberikan barisan eksak
% segment-id: o014.li2.chapter3.k-injectives.g004
	\[\begin{tikzcd}
		[column sep=small, every cell/.append style = {font = \small}]
		\displaystyle\varprojlim_k \displaystyle\prod_n
		\Hom\left( X^n, I_k^{n-1} \right) \arrow[r] &
		\displaystyle\varprojlim_k \displaystyle\prod_n
		\Hom\left( X^n, I_k^n \right) \arrow[r] &
		\displaystyle\varprojlim_k \displaystyle\prod_n
		\Hom\left( X^n, I_k^{n+1} \right)
	\end{tikzcd}\]
	Selanjutnya, pertukaran $\varprojlim_k$ dan $\prod_n$ memberikan barisan
	eksak\footnote{Catatan penerjemah (O014-C064): pada suku ketiga sumber,
	subskrip $k$ pada $\varprojlim_k$ hilang. Karena kalimat ini secara eksplisit
	menukar $\varprojlim_k$ dengan $\prod_n$, dan isomorfisme vertikal di bawahnya
	mengidentifikasi limit dengan $I^{n+1}$, target memulihkan subskrip tersebut.}
% segment-id: o014.li2.chapter3.k-injectives.g005
	\[\begin{tikzcd}
		[every cell/.append style = {font = \small}]
		\displaystyle\prod_n \varprojlim_k
		\Hom\left( X^n, I_k^{n-1} \right) \arrow[r] &
		\displaystyle\prod_n \varprojlim_k
		\Hom\left( X^n, I_k^n \right) \arrow[r] &
		\displaystyle\prod_n \varprojlim_k
		\Hom\left( X^n, I_k^{n+1} \right) \\
		\displaystyle\prod_n \Hom\left( X^n, I^{n-1} \right)
		\arrow[u, "\sim" sloped] &
		\displaystyle\prod_n \Hom\left( X^n, I^n \right)
		\arrow[u, "\sim" sloped] &
		\displaystyle\prod_n \Hom\left( X^n, I^{n+1} \right)
		\arrow[u, "\sim" sloped]
	\end{tikzcd}\]
	Ini juga merupakan penggalan dari $\Hom^\bullet(X,I)$, dan kohomologi
	pada suku tengahnya tepat sama dengan
	$\Hom_{\cate{K}(\mathcal{A})}(X,I)$. Dengan demikian,
	$\Hom_{\cate{K}(\mathcal{A})}(X,I)=0$. Terapkan Lema
	\ref{prop:K-injective-criterion} untuk menyelesaikan pembuktian.
\end{bukti}

% segment-id: o014.li2.chapter3.k-injectives.le003
\begin{lema}\label{prop:K-injective-truncated}
	Andaikan $\mathcal{A}$ mempunyai cukup banyak objek injektif. Untuk setiap
	$A\in\Obj(\cate{C}(\mathcal{A}))$, terdapat diagram komutatif dalam
	$\cate{C}(\mathcal{A})$
% segment-id: o014.li2.chapter3.k-injectives.g006
	\[\begin{tikzcd}
		\cdots \arrow[r] & I_2 \arrow[r] & I_1 \arrow[r] & I_0 \\
		\cdots \arrow[r] & \tau^{\geq -2} A\arrow[u, "f_2"'] \arrow[r] &
		\tau^{\geq -1} A \arrow[r] \arrow[u, "f_1"'] &
		\tau^{\geq 0} A \arrow[u, "f_0"']
	\end{tikzcd}\]
	dengan $\tau^{\geq k}A$ beserta panah-panah di antaranya seperti dalam
	Definisi \ref{def:truncation-cplx}, sedemikian sehingga:
% segment-id: o014.li2.chapter3.k-injectives.l006
	\begin{itemize}
% segment-id: o014.li2.chapter3.k-injectives.i012
		\item setiap $I_k$ merupakan objek
		$\cate{C}^+(\mathcal{A})$ yang tersusun atas objek-objek injektif;
% segment-id: o014.li2.chapter3.k-injectives.i013
		\item setiap panah vertikal $f_k$ merupakan monomorfisme sekaligus
		kuasi-isomorfisme;
% segment-id: o014.li2.chapter3.k-injectives.i014
		\item untuk setiap $k$ dan $n$, morfisme
		$I_{k+1}^n\to I_k^n$ merupakan epimorfisme yang mempunyai penampang.
	\end{itemize}
\end{lema}

% segment-id: o014.li2.chapter3.k-injectives.pf005
\begin{bukti}
	Kita membangun diagram dari kanan ke kiri. Semua resolusi injektif yang
	dibahas berikut ini secara implisit berupa monomorfisme.

	Pertama, terapkan Teorema \ref{prop:resolution-existence} pada
	$\tau^{\geq0}A\in\Obj(\cate{C}^+(\mathcal{A}))$ untuk memperoleh
	resolusi injektif $f_0:\tau^{\geq0}A\to I_0$.

	Selanjutnya, definisi langsung menunjukkan bahwa
	$\tau^{\geq-k-1}A\to\tau^{\geq-k}A$ merupakan epimorfisme pada setiap
	suku, sehingga merupakan epimorfisme di
	$\cate{C}^+(\mathcal{A})$. Andaikan kuasi-isomorfisme yang diperlukan
% segment-id: o014.li2.chapter3.k-injectives.q002
	\[
	\tau^{\geq0}A\xrightarrow{f_0}I_0,\quad\ldots,\quad
	\tau^{\geq-k}A\xrightarrow{f_k}I_k
	\]
	beserta diagram komutatif yang bersesuaian telah dibangun. Ambil resolusi
	injektif $H\xrightarrow{g}J$ dari
	$H:=\Ker\left[\tau^{\geq-k-1}A\to\tau^{\geq-k}A\right]$ melalui
	Teorema \ref{prop:resolution-existence}. Lalu, Proposisi
	\ref{prop:horseshoe} memberikan diagram komutatif dengan baris-baris eksak
	dalam $\cate{C}(\mathcal{A})$:
% segment-id: o014.li2.chapter3.k-injectives.g007
	\[\begin{tikzcd}
		0 \arrow[r] & J \arrow[r] & I_{k+1} \arrow[r] & I_k \arrow[r] & 0 \\
		0 \arrow[r] & H \arrow[r] \arrow[u, "g"] &
		\tau^{\geq-k-1}A \arrow[u, "{f_{k+1}}"] \arrow[r] &
		\tau^{\geq-k}A \arrow[r] \arrow[u, "f_k"'] & 0
	\end{tikzcd}\]
	sedemikian sehingga $f_{k+1}$ merupakan resolusi injektif. Keeksakan
	baris pertama setara dengan keeksakan
	$0\to J^n\to I_{k+1}^n\to I_k^n\to0$ untuk setiap $n$. Karena semua
	sukunya merupakan objek injektif di $\mathcal{A}$, Lema
	\ref{prop:inj-proj-split} menunjukkan bahwa
	$I_{k+1}^n\to I_k^n$ mempunyai penampang. Terbukti.
\end{bukti}

% segment-id: o014.li2.chapter3.k-injectives.p006
Pembahasan berikut menggunakan sejumlah notasi dari \S\ref{sec:lim1}.

% segment-id: o014.li2.chapter3.k-injectives.p007
Ambil $A$ seperti di atas. Ketika $k\geq0$ berubah, semua
$\tau^{\geq-k}A$ membentuk objek $\cate{InvSys}(\cate{C}(\mathcal{A}))$,
yang kita lambangkan dengan $\tau A$. Definisi funktor pemenggalan memberikan
keluarga morfisme kanonik $A\to\tau^{\geq-k}A$. Dengan meninjau kompleks
$\tau^{\geq-k}A$ suku demi suku, tampak bahwa $\varprojlim_k$-nya ada dan
morfisme-morfisme kanonik tersebut menginduksi isomorfisme
% segment-id: o014.li2.chapter3.k-injectives.e002
\begin{equation}\label{eqn:A-lim-truncation}
	A \rightiso \varprojlim_{k \geq 0} \left( \tau^{\geq -k} A \right).
\end{equation}

% segment-id: o014.li2.chapter3.k-injectives.p008
Sekarang tinjau $I_k$. Andaikan $\mathcal{A}$ mempunyai produk terhitung;
maka $\cate{C}(\mathcal{A})$ juga mempunyai produk terhitung yang dibangun
suku demi suku. Untuk $A\in\Obj(\cate{C}(\mathcal{A}))$, tinjau diagram
dalam Lema \ref{prop:K-injective-truncated}. Kita mempunyai
% segment-id: o014.li2.chapter3.k-injectives.q003
\[
I:=\varprojlim_k I_k \quad \text{ada}.
\]

% segment-id: o014.li2.chapter3.k-injectives.p009
Secara lebih konkret, berdasarkan Definisi \ref{def:Delta-diff} beserta
pembahasan sesudahnya, $I$ dan $A$ dapat dimasukkan ke dalam barisan-barisan
eksak
% segment-id: o014.li2.chapter3.k-injectives.q004
\begin{gather*}
	0 \to I \to \prod_k I_k \xrightarrow{\Delta_I} \prod_k I_k, \\
	0 \to A \to \prod_k \tau^{\geq -k} A
	\xrightarrow{\Delta_{\tau A}} \prod_k \tau^{\geq -k} A,
\end{gather*}
dengan $\Delta_I$ dan $\Delta_{\tau A}$ seperti dalam definisi tersebut.
Diperoleh diagram-diagram komutatif
% segment-id: o014.li2.chapter3.k-injectives.e003
\begin{equation}\label{eqn:K-injective-diagonal}\begin{tikzcd}
	I \arrow[hookrightarrow, r] \arrow[d] & \displaystyle\prod_k I_k \\
	\Cone(\Delta_I)[-1] \arrow[ru, "{\beta(\Delta_I)[-1]}"'] &
\end{tikzcd} \quad \begin{tikzcd}
	A \arrow[hookrightarrow, r] \arrow[d] &
	\displaystyle\prod_k \tau^{\geq -k} A \\
	\Cone(\Delta_{\tau A})[-1]
	\arrow[ru, "{\beta(\Delta_{\tau A})[-1]}"']
\end{tikzcd}\end{equation}
Kedua panah vertikal didefinisikan sebagai penyisipan ke koordinat pertama
dari kerucut pemetaan; pembaca dapat memeriksanya langsung. Penjelasan lain
yang lebih berat daripada yang diperlukan menggunakan sifat universal kernel
homotopi (Proposisi \ref{prop:homotopy-kernel-cokernel}).

% segment-id: o014.li2.chapter3.k-injectives.p010
Di sisi lain, dari $(f_k)_{k\geq0}$ dan sifat universal $\varprojlim$
diperoleh morfisme
% segment-id: o014.li2.chapter3.k-injectives.q005
\[
f:A\simeq\varprojlim_k(\tau^{\geq-k}A)\to\varprojlim_k I_k=I,
\]
yang membuat diagram berikut komutatif:
% segment-id: o014.li2.chapter3.k-injectives.e004
\begin{equation}\label{eqn:K-injective-A-I}\begin{tikzcd}
	I \arrow[r, "\text{seperti pada \eqref{eqn:K-injective-diagonal}}"
	inner sep=0.4em] & \Cone\left(\Delta_I\right)[-1] \\
	A \arrow[r] \arrow[u, "f"] &
	\Cone\left( \Delta_{\tau A} \right)[-1]
	\arrow[u, "\text{diinduksi oleh semua $f_k$}"'].
\end{tikzcd}\end{equation}
Lema \ref{prop:K-injective-projlimit-aux} bersama Lema
\ref{prop:K-injective-truncated} menunjukkan bahwa $I$ merupakan kompleks
K-injektif. Jika $f$ merupakan kuasi-isomorfisme, maka $f$ memberikan
resolusi K-injektif dari $A$. Langkah terakhir ini memerlukan syarat tambahan.

% segment-id: o014.li2.chapter3.k-injectives.le004
\begin{lema}\label{prop:K-injective-equiv}
	Dalam situasi di atas, andaikan lebih lanjut bahwa $\mathcal{A}$ mempunyai
	produk terhitung eksak (Konvensi \ref{con:exact-product}). Maka $f$
	merupakan kuasi-isomorfisme jika dan hanya jika
	$A\to\Cone(\Delta_{\tau A})[-1]$ merupakan kuasi-isomorfisme.
\end{lema}

% segment-id: o014.li2.chapter3.k-injectives.pf006
\begin{bukti}
	Tinjau diagram komutatif alami\footnote{Catatan penerjemah (O014-C065):
	panah horizontal bawah pertama pada sumber berlabel $\Delta_A$. Baris bawah
	merupakan sistem invers $\tau A$, dan ruas kanan pada baris yang sama sudah
	menggunakan $\alpha(\Delta_{\tau A})$ serta $\Cone(\Delta_{\tau A})$; target
	karena itu memulihkan label $\Delta_{\tau A}$.}
% segment-id: o014.li2.chapter3.k-injectives.g008
	\[\begin{tikzcd}
		\prod_k I_k \arrow[r, "\Delta_I"] &
		\prod_k I_k \arrow[r, "{\alpha(\Delta_I)}"] &
		\Cone\left(\Delta_I\right) \\
		\prod_k \tau^{\geq-k}A \arrow[r, "\Delta_{\tau A}"']
		\arrow[u, "{\prod_k f_k}"] &
		\prod_k \tau^{\geq-k}A \arrow[u, "{\prod_k f_k}"]
		\arrow[r, "{\alpha(\Delta_{\tau A})}"'] &
		\Cone\left(\Delta_{\tau A}\right) \arrow[u]
	\end{tikzcd}\]
	Dengan mengambil barisan eksak panjang kerucut pemetaan
	\eqref{eqn:cone-long-exact-sequence}, diperoleh diagram komutatif dengan
	baris-baris eksak:
% segment-id: o014.li2.chapter3.k-injectives.g009
	\[\begin{tikzcd}
		[column sep=tiny, every cell/.append style = {font = \small}]
		\cdots \Hm^n\left( \displaystyle\prod_k I_k \right) \arrow[r] &
		\Hm^n\left( \Cone(\Delta_I) \right) \arrow[r] &
		\Hm^{n+1}\left( \displaystyle\prod_k I_k \right)
		\arrow[r, "{\Hm^{n+1}(\Delta_I)}" inner sep=1em] &
		\Hm^{n+1}\left( \displaystyle\prod_k I_k \right) \\
		\cdots \Hm^n\left( \displaystyle\prod_k \tau^{\geq-k} A \right)
		\arrow[r] \arrow[u, "{\Hm^n(\prod_k f_k)}"] &
		\Hm^n\left( \Cone(\Delta_{\tau A}) \right)
		\arrow[r] \arrow[u] &
		\Hm^{n+1}\left(\displaystyle\prod_k \tau^{\geq-k} A \right)
		\arrow[u, "{\Hm^{n+1}(\prod_k f_k)}"]
		\arrow[r, "{\Hm^{n+1}(\Delta_{\tau A})}"' inner sep=1.5em] &
		\Hm^{n+1}\left( \displaystyle\prod_k \tau^{\geq-k} A \right)
		\arrow[u, "{\Hm^{n+1}(\prod_k f_k)}"'] .
	\end{tikzcd}\]
	Setiap $f_k$ merupakan kuasi-isomorfisme, dan keeksakan produk terhitung
	membuat $\prod_k f_k$ tetap menjadi kuasi-isomorfisme. Dengan menerapkan
	Proposisi \ref{prop:5-lemma}, diperoleh bahwa
	$\Cone(\Delta_{\tau A})\to\Cone(\Delta_I)$\footnote{Catatan penerjemah
	(O014-C066): sumber menulis $\Delta_{\tau_A}$ pada kemunculan ini, padahal
	objek sistem invers yang telah didefinisikan adalah $\tau A$ dan seluruh
	notasi di sekitarnya menggunakan $\Delta_{\tau A}$. Target menyeragamkannya
	dengan notasi yang telah didefinisikan.} juga merupakan
	kuasi-isomorfisme.

	Selanjutnya, kita buktikan bahwa $I\to\Cone(\Delta_I)[-1]$ merupakan
	kuasi-isomorfisme. Pertama, $\Delta_I$ merupakan epimorfisme; ini dapat
	diperiksa suku demi suku. Memang,
	$\Delta_{I^n}:\prod_k I_k^n\to\prod_k I_k^n$ merupakan epimorfisme karena
	$I_{k+1}^n\to I_k^n$ mempunyai penampang; lihat Contoh
	\ref{eg:lim1-section}. Bandingkan barisan eksak pendek
	$0\to I\to\prod_k I_k\xrightarrow{\Delta_I}\prod_k I_k\to0$
	dengan Lema \ref{prop:cone-qis}. Terlihat bahwa morfisme
	$I\to\Cone(\Delta_I)[-1]$ yang didefinisikan sebelumnya memang merupakan
	kuasi-isomorfisme.

	Masukkan hasil-hasil ini ke dalam diagram komutatif
	\eqref{eqn:K-injective-A-I}, lalu ambil kohomologi. Diperoleh syarat perlu
	dan cukup yang dikehendaki.
\end{bukti}

% segment-id: o014.li2.chapter3.k-injectives.th002
\begin{teorema}\label{prop:K-injective-resolution}
	Kategori abelian $\mathcal{A}$ yang memenuhi kedua syarat berikut mempunyai
	cukup banyak kompleks K-injektif:
% segment-id: o014.li2.chapter3.k-injectives.l007
	\begin{compactitem}
% segment-id: o014.li2.chapter3.k-injectives.i015
		\item $\mathcal{A}$ mempunyai produk terhitung eksak;
% segment-id: o014.li2.chapter3.k-injectives.i016
		\item $\mathcal{A}$ mempunyai cukup banyak objek injektif.
	\end{compactitem}
\end{teorema}

% segment-id: o014.li2.chapter3.k-injectives.pf007
\begin{bukti}
	Lema \ref{prop:K-injective-equiv} mereduksi masalah menjadi pembuktian
	bahwa $A\to\Cone(\Delta_{\tau A})[-1]$ merupakan kuasi-isomorfisme untuk
	setiap $A\in\Obj(\cate{C}(\mathcal{A}))$. Tetapkan $n\in\Z$. Karena
	produk terhitung eksak, Lema \ref{prop:truncation-ses} memberikan
	isomorfisme kanonik
% segment-id: o014.li2.chapter3.k-injectives.q006
	\[
	\Hm^n\left( \prod_k \tau^{\geq -k} A \right)
	\simeq \prod_k \Hm^n\left( \tau^{\geq -k} A \right)
	\simeq \prod_{k \geq -n} \Hm^n(A).
	\]
	Jika objek-objek $\mathcal{A}$ dapat dibahas melalui elemen, maka setelah
	mengambil $\Hm^n$, morfisme $\Delta_{\tau A}$ berbentuk
% segment-id: o014.li2.chapter3.k-injectives.g010
	\[\begin{tikzcd}[row sep=tiny]
		\prod_{k \geq -n} \Hm^n(A) \arrow[r, "\nabla^n"] &
		\prod_{k \geq -n} \Hm^n(A) \\
		(a_k)_{k \geq -n} \arrow[mapsto, r] &
		(a_{k+1} - a_k)_{k \geq -n}.
	\end{tikzcd}\]
	Jelas bahwa $\nabla^n$ merupakan epimorfisme (sebuah penampang dapat
	dituliskan secara eksplisit), sedangkan
	$\Ker(\nabla^n)=\left\{(a)_{k\geq-n}:a\in\Hm^n(A)\right\}$, yakni bagian
	diagonal. Pengamatan ini mudah dipindahkan ke kategori $\mathcal{A}$
	umum. Sekarang tinjau barisan eksak panjang kerucut pemetaan
% segment-id: o014.li2.chapter3.k-injectives.q007
	\[
	\cdots \to \Hm^n \Cone(\Delta_{\tau A})[-1]
	\to \prod_{k \geq -n} \Hm^n(A)
	\xrightarrow[\text{epi}]{\nabla^n}
	\prod_{k \geq -n} \Hm^n(A)
	\to \Hm^n \Cone(\Delta_{\tau A}) \cdots
	\]
	Barisan ini mengidentifikasi morfisme
	$\Hm^n\Cone(\Delta_{\tau A})[-1]\to\prod_{k\geq-n}\Hm^n(A)$ dengan
	penyisipan $\Ker(\nabla^n)$. Selain itu,
	\eqref{eqn:K-injective-diagonal} mengakibatkan diagram komutatif
% segment-id: o014.li2.chapter3.k-injectives.g011
	\[\begin{tikzcd}
		\Hm^n(A) \arrow[d]
		\arrow[hookrightarrow, r, "\text{morfisme diagonal}"] &
		\prod_{k \geq -n} \Hm^n(A) \\
		\Hm^n {\Cone(\Delta_{\tau A})[-1]} \arrow[hookrightarrow, ru] &
	\end{tikzcd}\]
	Dengan demikian, untuk setiap $n$, morfisme
	$\Hm^n(A)\to\Hm^n\left(\Cone(\Delta_{\tau A})[-1]\right)$ merupakan
	isomorfisme.
\end{bukti}

% segment-id: o014.li2.chapter3.k-injectives.p011
Dengan membalikkan panah melalui Proposisi
\ref{prop:K-sigma-duality}, diperoleh versi dual Teorema
\ref{prop:K-injective-resolution}.

% segment-id: o014.li2.chapter3.k-injectives.th003
\begin{teorema}\label{prop:K-projective-resolution}
	Kategori abelian $\mathcal{A}$ yang memenuhi kedua syarat berikut mempunyai
	cukup banyak kompleks K-projektif:
% segment-id: o014.li2.chapter3.k-injectives.l008
	\begin{compactitem}
% segment-id: o014.li2.chapter3.k-injectives.i017
		\item $\mathcal{A}$ mempunyai koproduk terhitung eksak;
% segment-id: o014.li2.chapter3.k-injectives.i018
		\item $\mathcal{A}$ mempunyai cukup banyak objek projektif.
	\end{compactitem}
\end{teorema}

% segment-id: o014.li2.chapter3.k-injectives.co001
\begin{korolari}\label{prop:Grothendieck-K-projective}
	Jika $\mathcal{A}$ merupakan kategori Grothendieck yang mempunyai cukup
	banyak objek projektif (Definisi \ref{def:Grothendieck-cat}), maka
	$\mathcal{A}$ juga mempunyai cukup banyak kompleks K-projektif.
\end{korolari}

% segment-id: o014.li2.chapter3.k-injectives.pf008
\begin{bukti}
	Menurut definisi, kategori Grothendieck $\mathcal{A}$ mempunyai koproduk
	terhitung. Proposisi \ref{prop:Grothendieck-cat-coprod-exact} menyatakan
	bahwa koproduk terhitung tersebut eksak. Terapkan Teorema
	\ref{prop:K-projective-resolution}.
\end{bukti}

% segment-id: o014.li2.chapter3.k-injectives.ex003
\begin{contoh}\label{eg:Mod-K-injectives}
	Misalkan $R$ gelanggang. Teorema \ref{prop:K-injective-resolution} dan
	Korolari \ref{prop:Grothendieck-K-projective} menunjukkan bahwa
	$R\dcate{Mod}$ mempunyai cukup banyak kompleks K-injektif dan
	K-projektif.
\end{contoh}

% segment-id: o014.li2.chapter3.k-injectives.p012
Sesungguhnya, setiap kompleks di atas kategori Grothendieck mempunyai resolusi
K-injektif, dan resolusi K-injektif tersebut dapat dibuat funktorial seperti
pada Teorema \ref{prop:Grothendieck-injective}. Inilah teorema utama
\cite{Serp03}; lihat juga \cite[Tag 079P]{stacks}. Hasil itu mencakup
sejumlah situasi yang lazim dalam geometri tetapi tidak tercakup oleh Teorema
\ref{prop:K-injective-resolution}.
